\documentclass[../../main.tex]{subfiles}
\begin{document}
\section{向量丛的局部截面与整体截面}\label{section:ch10-局部与整体截面}

设 $\pi\colon E\to M$ 为拓扑空间 $M$ 上秩 $k$ 向量丛。$E$ 在开集 $U\subseteq M$ 上的一个\textbf{截面}是一个连续映射 $\sigma\colon U\to E$，满足 $\pi\circ\sigma=\mathrm{Id}_U$。换言之，$\sigma$ 为每个点 $p\in U$ 选取 $E_p$ 中的一个向量 $\sigma(p)$（\cref{figure:ch10-局部截面}）。若 $U=M$，则 $\sigma$ 称为\textbf{整体截面}。一个\textbf{局部截面}是 $U$ 为某点邻域（而不一定是 $M$ 本身）的截面；若 $U$ 未明确给出，则“截面”指 $U\subseteq M$ 的任意开子集上的截面，若不加限定也可为任意（未必连续的）映射 $\sigma\colon U\to E$ 满足 $\pi\circ\sigma=\mathrm{Id}_U$。“截面”一词若不加进一步限定，总是指连续截面。

\begin{figure}[htbp]
\centering
\includegraphics[width=0.8\textwidth]{fig-10-3.png}
\caption{向量丛的一个局部截面}\label{figure:ch10-局部截面}
\end{figure}

$E$ 的\textbf{零截面}是整体截面 $\zeta\colon M\to E$，定义为
\begin{align*}
\zeta(p)=0\in E_p\quad\text{对每个 }p\in M.
\end{align*}

与向量场的情形一样，截面 $\sigma$ 的\textbf{支撑}是集合 $\{p\in M:\sigma(p)\neq0\}$ 的闭包。

\begin{exercise}\label{exercise:ch10-09}
证明每个向量丛的零截面是连续的，且每个光滑向量丛的零截面是光滑的。（提示：考虑 $\Phi\circ\zeta$，其中 $\Phi$ 是一个局部平凡化。）
\end{exercise}

\begin{example}[向量丛的截面]\label{example:ch10-向量丛的截面}
设 $M$ 为带边或不带边的光滑流形。
\begin{enumerate}[(a)]
\item $TM$ 的截面就是 $M$ 上的向量场。
\item 给定一个带边或不带边的浸入子流形 $S\subseteq M$，环境切丛 $TM|_S\to S$ 的一个截面称为\textbf{沿 $S$ 的向量场}。它是一个连续映射 $X\colon S\to TM$，满足对每个 $p\in S$ 有 $X_p\in T_pM$。这不同于\textbf{在} $S$ 上的向量场，后者在每点满足 $X_p\in T_pS$。
\item 若 $E=M\times\mathbb{R}^k$ 是乘积丛，则 $E$ 的截面与从 $M$ 到 $\mathbb{R}^k$ 的连续函数之间有自然的一一对应：一个连续函数 $F\colon M\to\mathbb{R}^k$ 决定了一个截面 $\widetilde F\colon M\to M\times\mathbb{R}^k$，由 $\widetilde F(x)=(x,F(x))$ 给出，反之亦然。若 $M$ 是带边或不带边的光滑流形，则截面 $\widetilde F$ 光滑当且仅当 $F$ 光滑。
\item 上一段中的对应给出了 $C^\infty(M)$ 与平凡线丛 $M\times\mathbb{R}\to M$ 的光滑截面空间之间的自然等同。
\end{enumerate}
\end{example}

设 $E\to M$ 为光滑向量丛。我们记 $\Gamma(E)$ 为 $E$ 的所有光滑整体截面的集合（为与切丛类比，光滑截面有时也称为\textbf{向量丛场}）。当 $E=TM$ 时，此空间记为 $\mathfrak X(M)$ 或 $\mathfrak X(M,\mathbb{R})$，即 $M$ 上光滑向量场的空间，与我们先前的记号一致。类似地，$M$ 上光滑函数的空间 $C^\infty(M)$ 常记为 $\Gamma(T^0M)$，其中 $T^0M=M\times\mathbb{R}$ 是 $M$ 上平凡的秩 $1$ 向量丛。在任一情形中，$\Gamma(E)$ 都是实数上的向量空间，也是 $C^\infty(M)$ 上的模。

$E$ 在开集 $U\subseteq M$ 上的一个\textbf{局部标架}是 $E$ 在 $U$ 上的 $k$ 个光滑截面 $\sigma_1,\dots,\sigma_k$，使得对每个 $p\in U$，$\sigma_1(p),\dots,\sigma_k(p)$ 构成 $E_p$ 的一组基。

\begin{example}[局部标架]\label{example:ch10-局部标架}
设 $M$ 为带边或不带边的光滑流形。
\begin{enumerate}[(a)]
\item 若 $E\to M$ 是向量丛，且 $\Phi\colon\pi^{-1}(U)\to U\times\mathbb{R}^k$ 是 $E$ 在 $U$ 上的局部平凡化，则由 $\sigma_i(p)=\Phi^{-1}(p,e_i)$ 定义的截面 $\sigma_1,\dots,\sigma_k$（其中 $e_1,\dots,e_k$ 是 $\mathbb{R}^k$ 的标准基）构成 $E$ 在 $U$ 上的一个光滑局部标架。特别地，每个光滑局部平凡化都决定一个光滑局部标架，反之亦然。
\item $TM$ 在 $M$ 的任一光滑坐标图 $(U,\varphi)$ 上的局部标架是坐标标架 $\partial/\partial x^1,\dots,\partial/\partial x^n$。
\item 若 $S\subseteq M$ 是带边或不带边的浸入子流形，则沿 $S$ 的环境切丛 $TM|_S$ 的一个光滑局部标架在每个 $p\in S$ 处给出切空间 $T_pM$ 的一组基，而不一定是 $T_pS$ 的基。例如，设 $M=\mathbb{R}^2$ 而 $S=\mathbb{S}^1$，则向量场 $\partial/\partial x$ 与 $\partial/\partial y$ 在 $\mathbb{S}^1$ 的任意点附近构成 $TM|_{\mathbb{S}^1}$ 的局部标架，尽管其中没有一个是切于 $\mathbb{S}^1$ 的。
\end{enumerate}
\end{example}

下面的命题推广了我们对向量场所知的几个事实，即\cref{proposition:ch08-11}~与\cref{proposition:ch08-14}。其证明要求你在习题~\ref{exercise:ch10-15}~中补出，它基本重复着早先对向量场所用的论证。

\begin{proposition}[光滑截面的存在性]\label{proposition:ch10-光滑截面存在性}
设 $E\to M$ 为带边或不带边光滑流形 $M$ 上的光滑向量丛。
\begin{enumerate}[(a)]
\item 对任意 $p\in M$，每个向量 $\tilde v\in E_p$ 都被包含在某个光滑整体截面中。
\item $E$ 的一个截面 $\sigma\colon M\to E$ 是光滑的，当且仅当 $M$ 的每一点都有邻域 $U$，使得 $\sigma|_U$ 是光滑的。
\item 若 $\sigma\colon U\to E$ 是开集 $U\subseteq M$ 上的光滑截面，且 $p\in U$，则存在一个光滑整体截面 $\widetilde\sigma\in\Gamma(E)$，在 $p$ 的某个邻域上等于 $\sigma$。
\end{enumerate}
\end{proposition}

关键工具是下面的引理，它推广了关于光滑向量场的相应结果（\cref{lemma:ch08-13}）。

\begin{lemma}[截面扩张引理]\label{lemma:ch10-截面扩张引理}
设 $E\to M$ 为光滑向量丛。若 $A\subseteq M$ 是闭集，而 $\sigma\colon A\to E$ 是连续截面（即对每个 $p\in A$ 有 $\sigma(p)\in E_p$），且在每个 $p\in A$ 处存在定义于 $p$ 的某个邻域上的光滑截面，其在 $A$ 上的限制等于 $\sigma$，则 $\sigma$ 可被扩张为一个光滑整体截面。
\end{lemma}

尽管我们对光滑向量场的许多构造都可平移到任意光滑向量丛上，但有一处不行：不存在切丛的整体光滑标架，因为光滑 $1$ 形式的整体存在性与流形的整体拓扑密切相关。（例如，我们将在第~15~章中证明奇数维紧流形总具有一个处处非零的向量场，而偶数维球面 $\mathbb{S}^{2n}$ 则没有。）

下面的命题是\cref{proposition:ch08-1}~的向量丛类比：它使我们能够把任意光滑局部线性无关截面组扩充为光滑局部标架。它比切丛情形中的相应结果（\cref{proposition:ch08-1}）稍强；这是它在习题~\ref{exercise:ch10-16}~中被要求证明的另一个原因。

\begin{proposition}[扩充为局部标架]\label{proposition:ch10-扩充为局部标架}
设 $E\to M$ 为光滑 $k$ 秩向量丛。若 $\sigma_1,\dots,\sigma_m$ 是开集 $U\subseteq M$ 上的光滑截面，且 $\sigma_1|_p,\dots,\sigma_m|_p$ 对某点 $p\in U$ 线性无关，则存在 $p$ 的一个邻域 $V\subseteq U$，使得 $\sigma_1,\dots,\sigma_m$ 在 $V$ 上可被扩充为 $E$ 的光滑局部标架。
\end{proposition}

关于光滑向量场，我们所证明的另一个重要事实是全局光滑向量场存在（\cref{proposition:ch08-14}），这是光滑向量丛的一般性质。

\begin{proposition}[整体光滑截面]\label{proposition:ch10-整体光滑截面}
设 $E\to M$ 为带边或不带边光滑流形 $M$ 上的光滑向量丛。给定任意点 $p\in M$ 与任意向量 $\tilde v\in E_p$，存在光滑整体截面 $\sigma\in\Gamma(E)$，使得 $\sigma(p)=\tilde v$。
\end{proposition}

\begin{proof}
选取 $E$ 在 $p$ 的邻域 $U$ 上的光滑局部平凡化 $\Phi\colon\pi^{-1}(U)\to U\times\mathbb{R}^k$。令 $v=(\pi_{\mathbb{R}^k}\circ\Phi)(\tilde v)\in\mathbb{R}^k$，并定义 $\widetilde\sigma_0\in\Gamma(\pi^{-1}(U))$ 为 $\widetilde\sigma_0(q)=\Phi^{-1}(q,v)$。它是光滑的，因为它是 $\Phi^{-1}$ 与 $q\mapsto(q,v)$ 的复合；并且根据 $v$ 的定义，$\widetilde\sigma_0(p)=\tilde v$。

令 $\psi\in C^\infty(M)$ 为 $p$ 的光滑 bump 函数，在 $p$ 的某邻域 $\widetilde U\subseteq U$ 上恒等于 $1$。定义 $\widetilde\sigma\colon M\to E$ 为
\begin{align*}
\widetilde\sigma(q)=
\begin{cases}
\psi(q)\,\widetilde\sigma_0(q),& q\in U,\\
0\in E_q,& q\notin\mathrm{supp}\,\psi.
\end{cases}
\end{align*}
容易验证，在 $U\setminus\mathrm{supp}\,\psi$ 与 $\widetilde U$ 的交上，这两个定义给出同一个截面，因此 $\widetilde\sigma$ 是良好定义的且在 $M$ 上连续。在 $U$ 上它是光滑的，而在 $M\setminus\mathrm{supp}\,\psi$ 上它恒为零，因而是光滑的；由于这两个集合覆盖 $M$，$\widetilde\sigma$ 是光滑的。最后，$\widetilde\sigma(p)=1\cdot\widetilde\sigma_0(p)=\tilde v$。
\end{proof}

下面的推论是光滑向量丛的单位分解定理。其证明要求你补出，完全模仿向量场相应结果的证明（习题~8-12）。

\begin{corollary}[向量丛的单位分解]\label{corollary:ch10-丛单位分解}
设 $M$ 为带边或不带边的光滑流形，$E\to M$ 为光滑向量丛。对 $M$ 的任意开覆盖 $\mathcal U$，存在 $\Gamma(E)$ 中的光滑整体截面族 $\{\psi_i\}$，满足：
\begin{enumerate}[(i)]
\item $\{\psi_i\}$ 的支撑形成 $M$ 上的局部有限族；
\item 对每个 $i$，存在某个 $U_j\in\mathcal U$ 使得 $\mathrm{supp}\,\psi_i\subseteq U_j$；
\item 对每个 $p\in M$，$\sum_i\psi_i(p)=\zeta(p)$，即 $E_p$ 中的零元（其中求和只有有限项非零）。
\end{enumerate}
\end{corollary}

最后，我们记录向量丛坐标的一种有用解释。你应该把它与切丛的坐标表示作比较。

\begin{example}[光滑向量丛的局部截面]\label{example:ch10-丛的局部截面}
设 $E\to M$ 为秩 $k$ 光滑向量丛，$(U,\Phi)$ 为 $E$ 在开集 $U\subseteq M$ 上的光滑局部平凡化。记 $\widehat\Phi=\varphi\times\mathrm{Id}_{\mathbb{R}^k}\circ\Phi\colon\pi^{-1}(U)\to\widehat U\times\mathbb{R}^k$，其中 $\widehat U=\varphi(U)\subseteq\mathbb{R}^n$。若 $S\subseteq E$ 是任意子集，则
\begin{align*}
S\cap\pi^{-1}(U)\text{ 是 }\pi^{-1}(U)\text{ 的嵌入子流形（带边或不带边）}
\end{align*}
当且仅当 $\widehat\Phi(S\cap\pi^{-1}(U))\subseteq\widehat U\times\mathbb{R}^k$ 是 $\widehat U\times\mathbb{R}^k$ 的嵌入子流形（带边或不带边）。因为 $\widehat\Phi$ 是微分同胚，这等价于说 $S$ 在 $E$ 中是嵌入子流形。
\end{example}

向量丛的光滑局部平凡化也可视为在 $M$ 上定义局部坐标的一种方式。若 $\Phi\colon\pi^{-1}(U)\to U\times\mathbb{R}^k$ 是这样的局部平凡化，而 $(U,\varphi)$ 是 $M$ 上的光滑坐标图，则 $\widehat\Phi=(\varphi\times\mathrm{Id}_{\mathbb{R}^k})\circ\Phi$ 把 $\pi^{-1}(U)$ 映到 $\mathbb{R}^{n+k}$ 中的开集 $\widehat U\times\mathbb{R}^k$ 上。若 $\varphi=(x^1,\dots,x^n)$，我们在 $\pi^{-1}(U)$ 上定义坐标 $(x^1,\dots,x^n,v^1,\dots,v^k)$，使得
\begin{align*}
\widehat\Phi(\tilde w)=(x^1(\pi(\tilde w)),\dots,x^n(\pi(\tilde w)),v^1(\tilde w),\dots,v^k(\tilde w)).
\end{align*}
这些称为由 $\Phi$ 决定的向量丛坐标。在这些坐标中，$\Phi$ 具有特别简单的形式
\begin{align*}
\Phi(x,v)=(x,v).
\end{align*}

由此可知，$\pi^{-1}(U)$ 上的光滑截面 $\sigma$ 在向量丛坐标下具有形式
\begin{align*}
\sigma(x)=\big(x,\sigma^1(x),\dots,\sigma^k(x)\big),
\end{align*}
其中 $\sigma^i(x)=v^i(\sigma(x))$ 是 $U$ 上的光滑函数；截面 $\sigma$ 光滑当且仅当这些分量函数是光滑的。更一般地，$E$ 在 $U$ 上的光滑局部标架 $\sigma_1,\dots,\sigma_k$ 可解释为 $U$ 上的一组向量值光滑函数 $\sigma_i=(\sigma_i^1,\dots,\sigma_i^k)$，使得矩阵 $\big(\sigma_i^j(x)\big)$ 对每个 $x\in U$ 可逆。
\begin{proposition}[向量丛坐标图]\label{proposition:ch10-向量丛坐标图}
设 $\pi\colon E\to M$ 为光滑向量丛，$(U,\Phi)$ 为光滑局部平凡化，而 $(\sigma_i)$ 为与之相应的光滑局部标架。若 $(U,\varphi)$ 是 $M$ 上的光滑坐标图，则
\begin{align*}
\widetilde\varphi=\big(\varphi\times\mathrm{Id}_{\mathbb{R}^k}\big)\circ\Phi
\end{align*}
给出 $\pi^{-1}(U)$ 上的光滑坐标图。
\end{proposition}

\begin{proof}
只需验证 $\widetilde\varphi=(\varphi\times\mathrm{Id}_{\mathbb{R}^k})\circ\Phi$，其中 $\Phi$ 是与 $(\sigma_i)$ 相联系的局部平凡化。作为微分同胚的复合，它是微分同胚。
\end{proof}

正如向量场的光滑性可用任一光滑坐标图中分量函数来刻画，向量丛截面的光滑性也可用局部标架来刻画。设 $(\sigma_i)$ 是 $E$ 在某开集 $U\subseteq M$ 上的光滑局部标架。若 $\tau\colon M\to E$ 是一个粗截面，则 $\tau$ 在任意点 $p\in U$ 处的值可唯一写为 $\tau(p)=\tau^i(p)\sigma_i(p)$。这定义了 $k$ 个函数 $\tau^i\colon U\to\mathbb{R}$，称为 $\tau$ 关于所给局部标架的\textbf{分量函数}。

\begin{proposition}[光滑性的局部标架判别法]\label{proposition:ch10-光滑性局部标架判别}
设 $\pi\colon E\to M$ 为光滑向量丛，$\tau\colon M\to E$ 为粗截面。若 $(\sigma_i)$ 是 $E$ 在开集 $U\subseteq M$ 上的光滑局部标架，则 $\tau$ 在 $U$ 上光滑，当且仅当其关于 $(\sigma_i)$ 的分量函数光滑。
\end{proposition}

\begin{proof}
令 $\Phi\colon\pi^{-1}(U)\to U\times\mathbb{R}^k$ 为与局部标架 $(\sigma_i)$ 相联系的局部平凡化。由于 $\Phi$ 是微分同胚，$\tau$ 在 $U$ 上光滑当且仅当复合映射 $\Phi\circ\tau$ 在 $U$ 上光滑。直接验证可得
\begin{align*}
\Phi\circ\tau(p)=\big(p,(\tau^1(p),\dots,\tau^k(p))\big),
\end{align*}
其中 $(\tau^i)$ 是 $\tau$ 关于 $(\sigma_i)$ 的分量函数，因此 $\Phi\circ\tau$ 光滑当且仅当分量函数 $\tau^i$ 光滑。
\end{proof}

\begin{exercise}\label{exercise:ch10-23}
设 $E\to M$ 为向量丛。证明 $E$ 的粗截面连续，当且仅当它在每个局部标架中的分量函数连续。
\end{exercise}

命题~\ref{proposition:ch10-光滑性局部标架判别}~同样适用于局部截面，因为 $E$ 在开集 $V\subseteq M$ 上的局部截面就是限制丛 $E|_V$ 的整体截面。

局部标架与局部平凡化之间的对应，导出下面这一刻画切丛光滑结构的唯一性结果。

\begin{proposition}[切丛光滑结构的唯一性]\label{proposition:ch10-切丛光滑结构唯一性}
设 $M$ 为带边或不带边的光滑 $n$ 流形。在\cref{proposition:ch03-18}~中构造的 $TM$ 上的拓扑与光滑结构，是使 $\pi\colon TM\to M$ 成为以所给纤维向量空间结构为结构的光滑向量丛、且所有坐标向量场都是光滑局部截面的唯一拓扑与光滑结构。
\end{proposition}

\begin{proof}
假设 $TM$ 被赋予了某个拓扑与光滑结构，使其成为具有所给性质的光滑向量丛。若 $(U,\varphi)$ 是 $M$ 的任意光滑坐标图，相应的坐标标架 $(\partial/\partial x^i)$ 是 $U$ 上的光滑局部标架，故由\cref{proposition:ch10-向量丛坐标图}~存在与此局部标架相联系的光滑局部平凡化 $\Phi\colon\pi^{-1}(U)\to U\times\mathbb{R}^n$。回顾\cref{example:ch10-局部标架}~的构造，可见此局部平凡化正是\cref{proposition:ch10-切丛作为向量丛}~中构造的映射 $\Phi$。由\cref{proposition:ch10-向量丛坐标图}，自然坐标图 $\widetilde\varphi=(\varphi\times\mathrm{Id}_{\mathbb{R}^n})\circ\Phi$ 属于所给光滑结构。于是所给光滑结构等于\cref{proposition:ch03-18}~中构造的那一个。
\end{proof}
\end{document}
